\UnnumberedChapter{KẾT LUẬN}

\textbf{Tổng kết.} Nhóm đã xây dựng quy trình học máy đầy đủ cho bài toán phân loại 43 loại biển báo giao thông trên bộ dữ liệu GTSRB: khám phá dữ liệu, tiền xử lý, so sánh ba cách biểu diễn (pixel, HOG, PCA) và hai thuật toán (K-NN, Cây quyết định), quét siêu tham số trên tập kiểm định chia theo track và đánh giá cuối trên tập Test. Các kết quả chính:

\begin{itemize}
    \item Mô hình tốt nhất là \textbf{K-NN kết hợp HOG} ($k=9$, metric euclidean) với Accuracy \textbf{0.9082} và Macro-F1 \textbf{0.8840} trên tập Test, so với baseline (pixel kết hợp K-NN) là 0.5253 và 0.4682.
    \item Biểu diễn đặc trưng quan trọng hơn việc chọn thuật toán: HOG vượt pixel ở cả K-NN và Cây quyết định.
    \item PCA (100 chiều) giữ gần như nguyên độ chính xác của HOG (0.9051) và rút ngắn thời gian dự đoán hơn một nửa.
    \item Lỗi tập trung ở các biển giới hạn tốc độ (khác nhau ở chữ số) và các biển cảnh báo hình tam giác có hình minh họa nhỏ, và nặng hơn với biển có kích thước ảnh nhỏ.
    \item Quy trình tái tạo được: hai lần chạy đầy đủ độc lập cho cùng Accuracy và Macro-F1 ở cả chín phương án.
\end{itemize}

\textbf{Hướng phát triển.}

\begin{itemize}
    \item Dùng thông tin màu (đặc trưng HSV, histogram màu) kết hợp HOG để phân biệt các biển cảnh báo hình tam giác.
    \item Tăng độ phân giải cho biển nhỏ hoặc trích đặc trưng riêng cho vùng chữ số của biển giới hạn tốc độ.
    \item Dùng đánh giá chéo theo track với nhiều phần để có khoảng tin cậy, và chọn mô hình chỉ theo tập kiểm định.
    \item Dùng chỉ mục láng giềng gần đúng (KD-tree, băm) để giảm thời gian dự đoán của K-NN.
    \item Phần mở rộng nếu có thêm thời gian: thử mạng nơ-ron tích chập hoặc mô hình huấn luyện sẵn để so sánh với các thuật toán đã học, kèm bộ phát hiện biển để không phụ thuộc vào ROI; xây dựng ứng dụng minh họa nhận ảnh và trả về nhãn dự đoán cùng độ tin cậy.
\end{itemize}
